
Since $ X$ is a random variable with a uniform distribution, by lemma 3.5, we have that 
for $ A \subset E$
\[
    \P_{X} (A) = \frac{|A|}{|E|} = \frac{|A|}{n} 
.\] 
The cumulative distribution function of the distribution of $ X$ is $\P_{X} (X \leq x)$, or more
specifically $ \P_{X} (e \in E: e \leq x)$
 
We see that clearly for $ x< 1$ $ F_{X}(x) = 0 $ since there are no elements of $ E$ that are
smaller
$ 1$. 

When $ x \in \R	$ is contained within $ [1,n)$, then there are at least some values of $ E$ that are 
smaller or equal to $ x$. Let's say that $ x \in [j, j+1)$, where $ j$ are natural numbers between
$ 1$ and $ n$. Let's see what the value of $ F_{X} $ is then in this case

\[
    F_{X} (x) = \P_{X} (e \in E : e \leq x) = \P_{X} (e \in E: e < j+1)  
.\] 

Using the fact that $ P_{X} $ is a uniform distribution, we get that 

\begin{align*}
    \P_{X} (e \in E: e < j+1) = \frac{| \{ e \in E: e < j+1 \} |}{n} = \frac{j}{n}  
\end{align*}

When $ x \in \R$ is greater or equal to $ n$ then all values of $ E$ are less than $ x$ meaning that 
\[
    \P_{X} (e \in E: e < n \leq x) = \frac{|e \in E: e < n|}{n} = \frac{n}{n} = 1
.\] 
